\documentclass[11pt,a4paper]{article}
\usepackage[T1]{fontenc}
\usepackage[utf8]{inputenc}
\usepackage{lmodern}
\usepackage{amsmath,amssymb,amsthm,mathtools}
\usepackage{geometry}
\geometry{margin=2.6cm}
\usepackage{booktabs,array}
\usepackage{enumitem}
\usepackage[dvipsnames]{xcolor}
\usepackage[hidelinks]{hyperref}
\usepackage{xurl}
\hypersetup{pdftitle={Three tiers of indeterminacy},pdfauthor={Toshihisa Urahigashi}}
\setlist{itemsep=2pt,topsep=4pt}
\raggedbottom
\emergencystretch=3em

\newcommand{\Z}{\mathbb{Z}}
\newcommand{\Q}{\mathbb{Q}}
\newcommand{\Zp}{\mathbb{Z}_p}
\newcommand{\Qp}{\mathbb{Q}_p}
\newcommand{\I}{\mathcal{I}}
\newcommand{\cO}{\mathcal{O}}
\newcommand{\fd}{\mathfrak{d}}
\newcommand{\tI}{(R_I)^{\sim}}
\newcommand{\hull}{\operatorname{hull}}
\newcommand{\mulog}{\mu^{\log}}
\newcommand{\Aut}{\operatorname{Aut}}
\newcommand{\Ism}{\operatorname{Ism}}
\newcommand{\Tr}{\operatorname{Tr}}
\newcommand{\GL}{\mathrm{GL}}
\newcommand{\ord}{\operatorname{ord}}
\newcommand{\qbb}{\underline{\underline{q}}}
\newcommand{\ceil}[1]{\lceil #1\rceil}
\newcommand{\floor}[1]{\lfloor #1\rfloor}
\newcommand{\cN}{\mathfrak{N}}
\newcommand{\cA}{\mathfrak{A}}
\newcommand{\cLf}{\mathfrak{L}_{\mathrm f}}
\newcommand{\cL}{\mathfrak{L}}
\newcommand{\leanus}{\textunderscore\allowbreak}
\newcommand{\lean}[1]{{\ttfamily\let\_\leanus #1}}

% labels
\newcommand{\lab}[2]{{\upshape\fcolorbox{#1}{#1!8}{\scriptsize\sffamily\bfseries\textcolor{#1}{#2}}}}
\newcommand{\tagS}{\lab{blue}{S}}
\newcommand{\tagX}{\lab{Purple}{X}}
\newcommand{\tagP}{\lab{ForestGreen}{P}}
\newcommand{\tagN}{\lab{BrickRed}{N}}
\newcommand{\tagH}{\lab{BurntOrange}{H}}

\theoremstyle{plain}
\newtheorem{theorem}{Theorem}[section]
\newtheorem{proposition}[theorem]{Proposition}
\newtheorem{corollary}[theorem]{Corollary}
\newtheorem{lemma}[theorem]{Lemma}
\theoremstyle{definition}
\newtheorem{definition}[theorem]{Definition}
\newtheorem{example}[theorem]{Example}
\theoremstyle{remark}
\newtheorem{remark}[theorem]{Remark}

\title{Three tiers of indeterminacy:\\ the exact hull at small places and the container of Step~(v) of IUT~IV}
\author{Toshihisa Urahigashi\\\small ORCID 0009-0004-0460-6242\\\small\texttt{123026.urahigashi.toshihisa@gmail.com}}
\date{Version 1.0.0, 8 November 2026\\\small\href{https://doi.org/10.5281/zenodo.24043000}{doi:10.5281/zenodo.24043000}}

\begin{document}
\maketitle

\begin{abstract}
In Step~(v) of the proof of inter-universal Teichm\"uller theory (IUT)~IV, Theorem~1.10, the holomorphic hull of
the union of the possible images of the $\Theta$-pilot object is bounded by a container that contains the region
$p^{\lambda}\cdot\tI$; the indeterminacies (Ind1) and (Ind2) are accounted for by the arbitrary automorphisms $\phi$ of
IUT~IV, Proposition~1.2, and (Ind3) by taking upper bounds. At places that are small in the sense of
Proposition~1.2(iii)---$p$ odd and every ramification index at most $p-2$---we compare this container with the hulls
reached by four classes of maps: unit scalars and field automorphisms; the maps induced factor by factor by the
actual (Ind1) and (Ind2); all factorwise automorphisms of the log-shells; and all automorphisms $\phi$ of
Proposition~1.2.

For the last two classes the hull is exactly $p^{\ceil{\lambda}-\fd_I}\cdot\tI$, where $\fd_I$ is the sum of the
orders of the differents, whereas the container of Step~(v) has hull $p^{\floor{\lambda}-\fd_I}\cdot\tI$. So when
$\lambda$ is not an integer---which is always the case at bad places---the container is larger than anything these
automorphisms reach by exactly a factor $p$, that is, by $\log p$ in log-volume. The difference comes from the proof
of Proposition~1.2(ii), which uses Proposition~1.1 with the different of every factor; Proposition~1.1 applied to the
factor that carries $p^{\lambda}$ removes it. Unit scalars and field automorphisms do not enlarge the hull. The actual
(Ind1), (Ind2) reach the exact value when $\lambda$ is an integer, provided that automorphisms of the local Galois
groups exist that carry a unit to the boundary of the log-shell; at small places we reduce this to one condition on
the transvections found by Hoshi--Nishio and Kondo. The difference is a term of conductor type and does not change the
shape of the inequality of IUT~IV.

No claim is made about the truth of Corollary~3.12 of IUT~III or of the $abc$ conjecture.
\end{abstract}

\noindent\textbf{MSC 2020:} 14H25 (primary); 11S15, 11S20.\quad
\textbf{Keywords:} inter-universal Teichm\"uller theory, indeterminacies, log-shell, holomorphic hull, different,
automorphisms of absolute Galois groups.

\section{Introduction}

\subsection{The question}
\tagS\ In Step~(v) of the proof of \cite[Thm.~1.10]{IUT4} (p.~27), Mochizuki writes:
\begin{quote}\small
``Thus, the \emph{inclusion} `$\phi(p^{\lambda}\cdot\tI)\subseteq p^{\floor{\lambda-\fd_I-a_I}}\cdot\log_p(R_I^{\times})$'
of Proposition 1.4, (iii), implies that the result of multiplying `$p^{\floor{\lambda}-|I|}\cdot2^{-|I|}\cdot
\log_p(R_I^{\times})$' by a suitable \emph{nonpositive} [\dots] integer power of $p_{v_{\Q}}$ \emph{contains} the
`union of possible images of a $\Theta$-pilot object' discussed in Step (iv). That is to say, the
\emph{indeterminacies} (Ind1) and (Ind2) are taken into account by the \emph{arbitrary} nature of the automorphism
`$\phi$' [cf. Proposition 1.2], while the \emph{indeterminacy} (Ind3) is taken into account by the fact that we are
considering \emph{upper bounds} [\dots]''
\end{quote}
So IUT~IV bounds the effect of (Ind1) and (Ind2) by all automorphisms that preserve the lattice of logarithms of
units. This note asks two questions at small places. How far do these automorphisms actually reach, and how does the
container compare with that? And how much of it can the maps induced by the actual (Ind1) and (Ind2) use?

\subsection{Results}
\begin{enumerate}[label=(\arabic*)]
\item \tagP\ For the class $\cL$ of all automorphisms $\phi$ of Proposition~1.2, and for the class $\cLf$ of
factorwise automorphisms of the log-shells, the hull of the union of the images of $p^{\lambda}\cdot\tI$ is exactly
$p^{\ceil{\lambda}-\fd_I}\cdot\tI$ (Theorem~\ref{t:exact}).
\item \tagP\ The container of Step~(v) has hull $p^{\floor{\lambda}-\fd_I}\cdot\tI$. When $\lambda\notin\Z$---always
the case at bad places---it is larger by exactly a factor $p$; when $\lambda\in\Z$ the two agree
(Corollary~\ref{c:compare}). The factor comes from the different of the factor that carries $p^{\lambda}$
(Remark~\ref{r:source}).
\item \tagP\ Unit scalars and field automorphisms do not enlarge the hull (Theorem~\ref{t:N}). The hull for the
actual (Ind1), (Ind2) is contained in the exact value of (1) (Theorem~\ref{t:Aupper}); when $\lambda\in\Z$ it is equal
to it, provided that at each place an automorphism of the local Galois group carries some integral element to the
boundary of the log-shell (Theorem~\ref{t:Areach}). At small places this condition follows from one condition on the
transvections of Hoshi--Nishio and Kondo (Proposition~\ref{p:transvection}); that condition is not proved here.
\item \tagP\ A positive room can occur only at $p=2$ or at primes with a ramified place above them, and only through
automorphisms of the local Galois groups that are not unit scalars times field automorphisms
(Proposition~\ref{p:where}).
\item \tagN\ For the explicit initial $\Theta$-datum D109 of \cite{Cat}, at its bad place above $109$, the container
exceeds the exact value by $\log 109$ for every label (Section~\ref{s:example}).
\end{enumerate}

\subsection{Relation to the author's earlier notes}
\tagP\ The author's note \cite{W311} treats a global, logical question: even granting the log-volume content of all
the results cited in the proof of Corollary~3.12, the conclusion does not follow. Its Theorem~22 builds a model on an
actual local field, replaces the groups of (Ind1), (Ind2) by all automorphisms of the log-shells (a choice it marks as
unproved), reads the possible images maximally, takes a container larger than that of IUT~IV, and uses only that the
enlargement of the hull is bounded by a constant $E_L$ depending on the local field. The present note is local and
numerical, and does not touch Step~(xi). It determines the exact value of that enlargement at small places, compares
it with the container of IUT~IV, and asks how much of the replaced group the actual (Ind1), (Ind2) reach. Its values
are smaller than the bound of \cite{W311}, so the conclusion of \cite{W311} is unaffected.

The author's note \cite{Gap} measures the (Ind2)-terms of lana-agents/iut from the region $\theta\cdot\bigotimes\I$
built on the log-shells; there the group of all factorwise automorphisms and the class of Proposition~1.2 give the
same value, and the span is the container of Step~(v) (Remark~\ref{r:frames}). The third report \cite{Third}
computes the floor slack $1+\floor{\lambda}-\lambda$ of Proposition~1.4(iii), and the fourth report \cite{Fourth}
records that $\Ism(G_K)=\Zp^{\times}$ for a finite extension $K$ of $\Qp$, after Minamide--Tsujimura \cite{MT}.

\subsection{Relation to earlier work}
Section~\ref{s:prior} gives details and quotations. In short: Proposition~1.1 of IUT~IV is stated with a free choice
of the factor that does not pay its different, but the proof of Proposition~1.2(ii) uses all factors; Dupuy--Hilado
note that one may choose the factor that makes the inclusion tightest, and exempt the factor with the largest different
\cite[Rem.~2.8.2, (6.4)]{DHPS}; lana-agents/iut states that for $\lambda=0$ the hull is exactly the container of
IUT~IV, and that the bound of IUT~IV adds $1$ for the rounding \cite{lana}. Mochizuki states that automorphisms of the
log-shells give only bounded discrepancies \cite[Ex.~3.5.1]{EssLgc}, and that the valuation is the data
``deformed'' by automorphisms that do not arise from field automorphisms \cite[\S I3]{AbsTopIII}. The exact value with
the ceiling for the region $p^{\lambda}\cdot\tI$, and the difference of exactly $\log p$ from the container when
$\lambda\notin\Z$, were not found in the literature we searched. The transvection formula of
Section~\ref{s:unitshear} follows in substance from Hoshi--Nishio \cite{HN} and Kondo \cite{Kondo}; it is not new.

\subsection{What is not claimed}
Nothing about the truth of Corollary~3.12 or of the $abc$ conjecture. That the family of possible images considered
here is the family of IUT: it does not include (Ind3), it takes the Kummer isomorphisms to be the identity, and it
treats label permutations as IUT~IV does (Section~\ref{s:classes}). How much room (Ind3) provides. Anything at places
that are not small. The effect on the final constants of Theorem~1.10; we do not compute it.

\subsection{Labels}
\tagS\ a statement about a cited text; \tagX\ a statement of another author; \tagP\ proved here; \tagN\ checked by an
exact computation (scripts in the deposit); \tagH\ not proved.

\section{Setting}\label{s:setting}

\subsection{Step (v)}
\tagS\ We use the notation of \cite[Prop.~1.1, 1.2, 1.4, Thm.~1.10, Steps (iv), (v)]{IUT4}. Let $p$ be an odd prime,
$j\in\{1,\dots,\ell^{*}\}$ a label and $I=\mathbb S^{\pm}_{j+1}=\{0,1,\dots,j\}$. For each $i\in I$ let $v_i$ be a
place of $K$ above $p$ (not necessarily distinct), $k_i=K_{v_i}$, $R_i=\cO_{k_i}$, $e_i$ its ramification index,
$\pi_i$ a uniformizer, and $\fd_i$ the order of a generator of the different of $R_i/\Zp$, with $\ord(p)=1$. Put
$R_I=\bigotimes_{\Zp}R_i$, let $\tI$ be its normalization, $T=R_I\otimes\Qp$, and $\fd_I=\sum_i\fd_i$. Following
Step~(v), put $i^{\dagger}=j$, let $\lambda=\ord(\qbb_{v_j}^{\,j^2})$ if $v_j$ is bad and $\lambda=0$ if it is good,
so that $\lambda\in\frac1{e_j}\Z$, and let $p^{\lambda}\in k_j$ have order $\lambda$. The region is
\[
\Theta=p^{\lambda}\cdot\tI\qquad(p^{\lambda}\text{ in the factor }j).
\]
For $x\in T$ written as $p^{k}\prod_i\pi_i^{a_i}$ ($k,a_i\in\Z$) the order $\ord(x)=k+\sum_ia_i/e_i$ is the same in
every component of $T$, and $\mulog(x\cdot\tI)-\mulog(\tI)=-\ord(x)\log p$ for the normalized log-volume of
\cite[Prop.~1.4]{IUT4}.

\subsection{Small places}
\begin{definition}\label{d:small}
A place $v\mid p$ is \emph{small} if $p$ is odd and $e_v\le p-2$ (as in \cite{Gap}); $p$ is small if every place
above it is small.
\end{definition}

\begin{lemma}\label{l:small}
\tagP\ If all factors are small, then:
\begin{enumerate}[label=(\alph*)]
\item each $k_i/\Qp$ is tamely ramified, $\fd_i=(e_i-1)/e_i$, $a_i=1/e_i$, hence $\fd_i+a_i=1$;
\item the log-shell is $\I_i=(2p)^{-1}\log_p(R_i^{\times})=\pi_i^{1-e_i}R_i$;
\item the module of the inclusion used in Step~(v) is
$p^{\floor{\lambda-\fd_I-a_I}}\cdot\log_p(R_I^{\times})=p^{\floor{\lambda}-|I|}\cdot2^{-|I|}\cdot
\log_p(R_I^{\times})=p^{\floor{\lambda}}\cdot\bigotimes_i\I_i$, so the ``suitable nonpositive integer power'' can be
taken to be $p^0$. We call it the \emph{container} and write $C$.
\end{enumerate}
\end{lemma}
\begin{proof}
(a) $e_i\le p-2<p$, so $p\nmid e_i$ and the different is $\pi_i^{e_i-1}R_i$; $\ceil{e_i/(p-2)}=1$. (b) \tagS\
\cite[Prop.~1.2(i)]{IUT4}: for $p>2$ and $e_i\le p-2$, $\log_p(R_i^{\times})=p^{a_i}\cdot R_i=\pi_iR_i$; since $p$
is odd, $2$ is a unit and $(2p)^{-1}\pi_iR_i=\pi_i^{1-e_i}R_i$. (c) By (a), $\fd_I+a_I=|I|$ is an integer.
\end{proof}

\subsection{Hull and room}
\tagS\ The holomorphic hull $\hull(X)$ of $X\subseteq T$ is the intersection of the sets $c\cdot\tI$, $c\in T^{\times}$,
that contain $X$ \cite[Rem.~3.9.5(i)]{IUT3}. For a set $\mathfrak X$ of $\Qp$-linear automorphisms of $T$ let
$U_{\mathfrak X}=\bigcup_{g\in\mathfrak X}g(\Theta)$ and define the \emph{room}
\[
\rho(\mathfrak X)=\mulog\bigl(\hull(U_{\mathfrak X})\bigr)-\mulog(\Theta);
\]
for a module $M\supseteq\Theta$ write $\rho(M)=\mulog(\hull(M))-\mulog(\Theta)$.

\subsection{Four classes}\label{s:classes}
\tagS\ In \cite[Thm.~3.11(i)]{IUT3} (p.~154), (Ind1) consists of ``the indeterminacies induced by the automorphisms
of the procession of $\mathcal D^{\vdash}$-prime-strips'', and (Ind2), at a nonarchimedean $v_{\Q}$, of ``the
indeterminacies induced by the action of independent copies of $\Ism$ [\dots] on each of the direct summands of the
$j+1$ factors appearing in the tensor product''. An automorphism of a $\mathcal D^{\vdash}$-prime-strip is, at each
place, an automorphism of the topological group $G_v$, which need not arise from a field automorphism.

\begin{definition}\label{d:classes}
Consider the following sets of $\Qp$-linear automorphisms of $T$:
\begin{itemize}
\item $\cN$: the maps $\bigotimes_ic_i\sigma_i$ with $c_i\in\Zp^{\times}$ and $\sigma_i\in\Aut(k_i/\Qp)$;
\item $\cA$: the maps $\bigotimes_ic_i\gamma_i$ with $c_i\in\Zp^{\times}$ and $\gamma_i$ the map induced on $k_i$, through
the mono-analytic reconstruction of the log-shell, by an automorphism $\beta_i$ of the topological group $G_{k_i}$;
\item $\cLf$: the maps $\bigotimes_ig_i$ with $g_i$ a $\Qp$-linear automorphism of $k_i$ such that $g_i(\I_i)=\I_i$;
\item $\cL$: the $\Qp$-linear automorphisms $\phi$ of $T$ with $\phi(\bigotimes_i\I_i)=\bigotimes_i\I_i$, that is,
the automorphisms of Proposition~1.2 of \cite{IUT4}.
\end{itemize}
\end{definition}

\begin{proposition}\label{p:incl}
\tagP\ $\cN\subseteq\cA\subseteq\cLf\subseteq\cL$.
\end{proposition}
\begin{proof}
A field automorphism of $k_i$ over $\Qp$ is induced by conjugation by an extension to an algebraic closure, and
\cite[Prop.~9]{Fourth} applies \cite[Thm.~F, Thm.~7.4]{MT}: $\Ism(G_{k_i})=\Zp^{\times}$, acting by unit scalars.
The mono-analytic log-shell is constructed from the $\mathcal D^{\vdash}$-prime-strip by a functorial algorithm
\cite[Prop.~1.2(vi)]{IUT3}, so the induced maps preserve the log-shells. Tensor products of lattice automorphisms
preserve the tensor product of the lattices.
\end{proof}

\begin{remark}[label permutations and (Ind3)]
(Ind1) also permutes the labels. As in Step~(v), we fix the factor $i^{\dagger}=j$ carrying the theta value; IUT~IV
symmetrizes over the choice of $i^{\dagger}$ when it averages \cite[p.~28]{IUT4}, and lana-agents/iut likewise takes
the hull for each (Ind1)-image separately \cite{lana}. We do not include (Ind3).
\end{remark}

\section{The exact hull for the classes of Proposition 1.2}\label{s:exact}

\begin{proposition}\label{p:inside}
\tagP\ If all factors are small, then $\Theta\subseteq p^{\ceil{\lambda}}\cdot\bigotimes_i\I_i$.
\end{proposition}
\begin{proof}
\tagS\ \cite[Prop.~1.1]{IUT4}: ``Fix an element $*\in I$; write $I^{*}=I\setminus\{*\}$. Then
$p^{\fd_{I^*}}\cdot\tI\subseteq R_I\subseteq\tI$,'' where $p^{\fd_{I^*}}$ is a product of generators $\delta_i\in R_i$
of the differents for $i\in I^{*}$. Take $*=j$. Then
$\Theta\subseteq p^{\lambda}R_j\otimes\bigotimes_{i\ne j}\delta_i^{-1}R_i$. By Lemma~\ref{l:small},
$\delta_i^{-1}R_i=\pi_i^{1-e_i}R_i=\I_i$. Finally $p^{\lambda}R_j\subseteq p^{k}\I_j=\pi_j^{e_jk+1-e_j}R_j$ if and only
if $k\le\lambda+(e_j-1)/e_j$; as $\lambda\in\frac1{e_j}\Z$, the largest such $k$ is $\ceil{\lambda}$.
\end{proof}

\begin{theorem}[exact hull]\label{t:exact}
\tagP\ Suppose all factors are small, and put $\mu=p^{\ceil{\lambda}}\prod_i\pi_i^{1-e_i}\in T^{\times}$. For
$\mathfrak X=\cLf$ and for $\mathfrak X=\cL$,
\[
\hull(U_{\mathfrak X})=\mu\cdot\tI=p^{\ceil{\lambda}-\fd_I}\cdot\tI,\qquad
\rho(\mathfrak X)=\bigl(\fd_I-(\ceil{\lambda}-\lambda)\bigr)\log p .
\]
\end{theorem}
\begin{proof}
Upper bound: every $\phi\in\cL$ preserves $p^{\ceil{\lambda}}\bigotimes_i\I_i=\mu R_I$, so
$U_{\cL}\subseteq\mu R_I$ by Proposition~\ref{p:inside}. Since $1\in R_I\subseteq\tI$, a set $c\cdot\tI$ containing
$\mu R_I$ contains $\mu$, hence $\mu\cdot\tI$; so $\hull(\mu R_I)=\mu\cdot\tI$ and $\hull(U_{\cL})\subseteq\mu\cdot\tI$.

Lower bound: the automorphism group of a free $\Zp$-module acts transitively on its primitive vectors. For $i\ne j$,
both $1$ and $\pi_i^{1-e_i}$ are primitive in $\I_i$ (they lie in $\I_i$ and not in $p\I_i=\pi_iR_i$), so some
$g_i\in\GL(\I_i)$ has $g_i(1)=\pi_i^{1-e_i}$. For $i=j$, both $p^{\lambda}$ and $p^{\ceil{\lambda}}\pi_j^{1-e_j}$ are
primitive in $p^{\ceil{\lambda}}\I_j$ (because $e_j\lambda\le e_j\ceil{\lambda}$), so some $g_j\in\GL(\I_j)$ has
$g_j(p^{\lambda})=p^{\ceil{\lambda}}\pi_j^{1-e_j}$. The image of $1\otimes\cdots\otimes1\otimes p^{\lambda}\in\Theta$
under $\bigotimes_ig_i\in\cLf$ is $\mu\in T^{\times}$, so $\hull(U_{\cLf})\supseteq\mu\cdot\tI$. With
$U_{\cLf}\subseteq U_{\cL}$ both hulls equal $\mu\cdot\tI$.

Room: $\ord(\mu)=\ceil{\lambda}-\fd_I$ by Lemma~\ref{l:small}(a), and $\ord(p^{\lambda})=\lambda$.
\end{proof}

\begin{corollary}[container and bound]\label{c:compare}
\tagP\ Under the hypotheses of Theorem~\ref{t:exact},
\begin{align*}
\hull(C)&=p^{\floor{\lambda}-\ceil{\lambda}}\cdot\hull(U_{\cL}),\\
\rho(C)-\rho(\cL)&=(\ceil{\lambda}-\floor{\lambda})\log p,\\
(\fd_I+1)\log p-\rho(\cL)&=(1+\ceil{\lambda}-\lambda)\log p .
\end{align*}
In particular, if $\lambda\notin\Z$ the hull of the container is $p^{-1}$ times the hull reached by all of $\cL$ (its
log-volume is larger by exactly $\log p$), and if $\lambda\in\Z$ they coincide. Here $(\fd_I+1)\log p$ is the room
allowed by \cite[Prop.~1.2(iii)]{IUT4}, $\phi(p^{\lambda}\cdot\tI)\subseteq p^{\lambda-\fd_I-1}\cdot\tI$.
\end{corollary}
\begin{proof}
By Lemma~\ref{l:small}(c), $C=p^{\floor{\lambda}}\bigotimes_i\I_i$, whose hull is
$p^{\floor{\lambda}}\prod_i\pi_i^{1-e_i}\cdot\tI$ as in the proof of Theorem~\ref{t:exact}. Compare orders.
\end{proof}

\begin{remark}[where the factor $p$ comes from]\label{r:source}
\tagS\ The proof of \cite[Prop.~1.2(ii)]{IUT4} (p.~11) reads: ``Observe that it follows from the first displayed
inclusion [\emph{of $R_I$-modules}!] of Proposition 1.1 that
$p^{\fd_I+a_I}\cdot\tI\subseteq\bigotimes_{i\in I}p^{a_i}\cdot R_i$''. \tagP\ Here the different of every factor is
paid. Using Proposition~1.1 with $*=i^{\dagger}$ instead gives, by the same argument,
\[
\phi(p^{\lambda}\cdot\tI)\subseteq p^{\floor{\lambda-\fd_{I\setminus\{i^{\dagger}\}}-a_I}}\cdot\log_p(R_I^{\times}).
\]
If all factors are small, $\fd_{I\setminus\{i^\dagger\}}+a_I=|I|-1+1/e_j$, and for $\lambda\in\frac1{e_j}\Z$ one has
$\floor{\lambda-1/e_j}=\ceil{\lambda}-1$, so the right-hand side is exactly $p^{\ceil{\lambda}}\bigotimes_i\I_i$, the
module of Proposition~\ref{p:inside}. So the factor $p$ of Corollary~\ref{c:compare} is the price of paying the
different of the factor that carries $p^{\lambda}$. The remaining $(1+\floor{\lambda}-\lambda)\log p$ between the
container and the bound of Proposition~1.2(iii) is the floor slack of Proposition~1.4(iii), computed on explicit data
in \cite[\S\,The floor]{Third}. None of this is a defect of IUT~IV: the estimate is an upper bound, and an upper bound
may be generous.
\end{remark}

\begin{remark}[the frame of \cite{Gap}]\label{r:frames}
\tagP\ \cite{Gap} measures from the region $\theta\cdot\bigotimes\I$ built on the log-shells. There the factorwise
group (\texttt{tensorAut} of lana-agents/iut) and the class of Proposition~1.2 give the same value, and the span of the
union is $p^{\kappa}\bigotimes\I$ with $\kappa=\floor{\ord_p\theta}$ at small places, which is the container $C$
\cite[Thm.~3.2, \S6.2]{Gap}. Here we measure from the region $\Theta=\theta\cdot\tI$ of Step~(v) itself. Again the
two classes agree (Theorem~\ref{t:exact}), but their common value is smaller than $C$ by a factor $p$ when
$\lambda\notin\Z$. The two statements are compatible because the starting regions differ:
$\rho(\theta\cdot\bigotimes\I)=\fd_I\log p$ exceeds $\rho(\cL)=(\fd_I-(\ceil{\lambda}-\lambda))\log p$.
\end{remark}

\begin{remark}[checks]\label{r:checks}
\tagN\ For $p=5$ and all factors equal to $k=\Q_5(\sqrt5)$ ($e=2\le p-2$), with two factors ($j=1$) and three
factors ($j=2$), we computed by exact linear algebra the largest $k$ with $\Theta\subseteq5^{k}\bigotimes\I$, for
$\theta=\sqrt5^{\,m}$ in the last factor: it is $\ceil{m/2}$ for $0\le m\le12$ (two factors) and $0\le m\le7$ (three
factors). For random elements of $\GL(\I)$ in each factor (400 and 150 samples per $m$ for two factors, 150 for three)
the componentwise minimal valuation of the images is $2\ceil{m/2}-(j+1)$ in $\sqrt5$-units, the value of
Theorem~\ref{t:exact}; the container gives $2\floor{m/2}-(j+1)$, which differs by $2$ exactly when $m$ is odd. For
three factors and $m=0$ the hull has radius $5^{3/2}$, in agreement with the statement of lana-agents/iut quoted in
Section~\ref{s:prior}.
\end{remark}

\section{The actual indeterminacies}\label{s:actual}

\begin{theorem}\label{t:N}
\tagP\ $\rho(\cN)=0$.
\end{theorem}
\begin{proof}
$\sigma_j(p^{\lambda})$ is a unit multiple of $p^{\lambda}$, and $c_i\sigma_i$ preserves $R_i$, hence the normalization;
so every element of $\cN$ maps $\Theta$ onto itself.
\end{proof}

\begin{theorem}\label{t:Aupper}
\tagP\ If all factors are small, then $\hull(U_{\cA})\subseteq p^{\ceil{\lambda}-\fd_I}\cdot\tI$, and
$\rho(\cA)\le(\fd_I-(\ceil{\lambda}-\lambda))\log p$.
\end{theorem}
\begin{proof}
Proposition~\ref{p:incl} and Theorem~\ref{t:exact}.
\end{proof}

\begin{definition}\label{d:U}
Say that a place $v$ satisfies \emph{(U)} if there are an automorphism $\beta$ of $G_{K_v}$ and $x\in\cO_{K_v}$ such
that the induced map $\gamma$ satisfies $|\gamma(x)|=|\pi_v|^{1-e_v}$, that is, $\gamma$ carries an integral element
to the boundary of the log-shell. (For $e_v=1$ the identity does.)
\end{definition}

\begin{theorem}\label{t:Areach}
\tagP\ Suppose all factors are small, $\lambda\in\Z$, and every $v_i$ satisfies (U). Then
$\hull(U_{\cA})=p^{\lambda-\fd_I}\cdot\tI$, the value of Theorem~\ref{t:exact}.
\end{theorem}
\begin{proof}
Take $x_i$, $\gamma_i$ as in (U) for $v_i$. Since $x_i$ is integral and $p^{\lambda}\in\Qp$,
$y=x_0\otimes\cdots\otimes x_{j-1}\otimes p^{\lambda}x_j\in\Theta$. Its image under $\bigotimes_i\gamma_i\in\cA$ is
$p^{\lambda}\,\gamma_0(x_0)\otimes\cdots\otimes\gamma_j(x_j)$, a pure tensor of nonzero elements, hence a unit of $T$
of order $\lambda+\sum_i(1-e_i)/e_i=\lambda-\fd_I$ in every component. So the hull contains
$p^{\lambda-\fd_I}\cdot\tI$; Theorem~\ref{t:Aupper} gives the other inclusion.
\end{proof}

\subsection{When does (U) hold?}\label{s:unitshear}
\tagX\ Automorphisms of $G_L$ that do not preserve the ring of integers of $\overline L$ exist
\cite[Rem.~4.3.1(ii)]{Hoshi}, and the group-theoretic log-shell equals the ring of integers exactly when $p$ is odd and
$k$ is unramified, or $k\cong\Q_2$ or $\Q_3(\zeta_3)$ \cite[Thm.~A(ii)]{HoshiTop}. Every such automorphism $\alpha$
induces a map $\alpha_+$ that is compatible with $\Tr_{k/\Qp}$ and restricts to an automorphism of its kernel
\cite[Lemma~2.3(ii)]{HN}. For $p$ odd and $[k:\Qp]$ even, the Jannsen--Wingberg presentation \cite{JW} as made
explicit by Kondo gives an automorphism $\phi$ (with $\phi(x_2)=x_2x_1$) whose induced map moves one generator of the
logarithms of units by a multiple of another, while the kernel of the trace is spanned by the others
\cite[Thm.~1.3, Lemma~1.5]{Kondo}; Hoshi--Nishio give unipotent induced maps $\alpha_+$ with
$(\alpha_+^n-\mathrm{id})^2=0$ \cite[Thm.~1.5]{HN}. Together these give induced maps of the form
\[
x\longmapsto x+\Tr_{k/\Qp}(x)\,\eta,\qquad\Tr_{k/\Qp}(\eta)=0 .
\]
This formula follows in substance from \cite{HN,Kondo}; we do not claim it.

\begin{proposition}\label{p:transvection}
\tagP\ Let $v$ be small with $e_v\ge2$. If some automorphism of $G_{K_v}$ induces $x\mapsto x+\Tr(x)\eta$ with
$|\eta|=|\pi_v|^{1-e_v}$, then $v$ satisfies (U).
\end{proposition}
\begin{proof}
$K_v/\Qp$ is tame, so $\Tr(\cO_{K_v})=\Zp$; take $x\in\cO_{K_v}$ with $\Tr(x)=1$. Then $\gamma(x)=x+\eta$ with
$|\eta|>1\ge|x|$, so $|\gamma(x)|=|\eta|$.
\end{proof}

\tagH\ Since the induced maps preserve the log-shell, $|\eta|\le|\pi_v|^{1-e_v}$ always; whether equality holds for
a suitable presentation is not decided here. It concerns the position of the units attached to the generators, and
does not follow from the statements we cite. For $[K_v:\Qp]$ odd, and for $\lambda\notin\Z$ (which needs a move by
$e_j\lambda\bmod e_j-1$ steps in the last factor rather than $e_j-1$), we have nothing.

\subsection{Where the room can come from}
\begin{proposition}\label{p:where}
\tagP\ (a) If $p$ is odd and every place above $p$ is unramified, then $\rho(\cL)=0$; so positive room needs $p=2$ or
a ramified place above $p$. (b) If $\rho(\cA)>0$, then some element of $\cA$ has a factor $c_i\gamma_i$ with
$\gamma_i$ not of the form $c\,\sigma$ ($c\in\Zp^{\times}$, $\sigma$ a field automorphism).
\end{proposition}
\begin{proof}
(a) \tagS\ \cite[Prop.~1.2(iv)]{IUT4}: if $p>2$ and $e_i=1$ for all $i$, then $\phi(\tI)\subseteq\tI$. Here
$\lambda\in\Z$, so $p^{\lambda}$ commutes with $\phi$. (b) Theorem~\ref{t:N}.
\end{proof}

\begin{remark}[the inequality of Step (xi)]
\tagP\ Write the inequality of the form of Step~(xi-f) for the family generated by $\cA$ (without (Ind3)) with the
normalizations of IUT~IV: the $q$-pilot has log-volume $-|\log q|$ and the native $\Theta$-pilot, averaged over the
labels, $-c\,|\log q|$ with $c=\ell(\ell+1)/12$. The inequality then reads $(c-1)|\log q|\le A+\sum_p\bar\rho_p$, where
$A$ is the archimedean part and $\bar\rho_p$ the room at $p$ averaged over the labels. When the height is so large that
$A<(c-1)|\log q|$, some $\bar\rho_p$ must be positive, so by Proposition~\ref{p:where} some actual automorphism of a
local Galois group that is not a unit scalar times a field automorphism must move the region, at $p=2$ or at a ramified
prime. This is a necessary condition only; it says nothing about whether the inequality holds. At small primes
$\bar\rho_p$ is at most the label average of the exact value of Theorem~\ref{t:exact}. Consistently with
\cite[p.~101]{EssLgc}, none of these terms involves the size of the $q$-orders.
\end{remark}

\section{An example}\label{s:example}
\tagN\ The explicit initial $\Theta$-datum D109 of \cite{Cat} ($\ell=7$, $\ell^{*}=3$) has a bad place above $p=109$
with $e=105\le107$ and $\ord_p(\qbb)=29/7$; the other place above $109$ is good, with $e=1$, so $109$ is small.
Take the component with every factor at this bad place. Then $\lambda=29j^2/7\notin\Z$ for $j=1,2,3$ (as $7\nmid29$)
and $\fd_I=104(j+1)/105$. In units of $\log109$:

\begin{center}
\begin{tabular}{rrrrrrr}
\toprule
$j$ & $\lambda$ & $\fd_I$ & exact (Thm.~\ref{t:exact}) & container & IV bound & container $-$ exact\\
\midrule
1 & 29/7 & 208/105 & 118/105 & 223/105 & 313/105 & 1\\
2 & 116/7 & 104/35 & 89/35 & 124/35 & 139/35 & 1\\
3 & 261/7 & 416/105 & 341/105 & 446/105 & 521/105 & 1\\
\bottomrule
\end{tabular}
\end{center}
Here ``IV bound'' is the room $(\fd_I+1)\log p$ allowed by Proposition~1.2(iii). With $\log109\approx4.691$ the three exact rooms are about $5.272$, $11.929$ and $15.236$. Which bad places are
selected for $\mathbb V^{\mathrm{bad}}_{\mathrm{mod}}$ does not affect this local computation.

\section{Discussion}
\begin{itemize}
\item \tagP\ The estimate of IUT~IV is an upper bound, and nothing in its logic is affected by the container being
larger than what is reached. Corollary~\ref{c:compare} states precisely how generous it is at small places. The
difference $(\ceil{\lambda}-\floor{\lambda})\log p$ does not involve the size of the $q$-orders: it is a term of
conductor type and does not change the shape of the height inequality.
\item \tagP\ At small places, enlarging the factorwise automorphisms $\cLf$---the group \texttt{tensorAut} of
lana-agents/iut---to all of $\cL$ does not change the hull. Any effect of the enlargement must come from places that
are not small.
\item \tagP\ The difference between the actual $\cA$ and $\cLf$ is a question about automorphisms of absolute Galois
groups that do not come from field automorphisms: how far they move integral structures. Unit scalars and field
automorphisms give no room (Theorem~\ref{t:N}).
\item \tagH\ We do not include (Ind3). IUT~IV accounts for it by considering upper bounds, and Mochizuki states that
(Ind3) by itself accounts for the possibility of confusing $q$ with $q^{j^2}$ \cite[p.~101]{EssLgc}. Where the family of
possible images with (Ind3) lies between the exact value and the container is not decided here.
\item \tagP\ The bound $E_L$ of \cite[Thm.~22(a)]{W311} reads, at small places where $b_I+a_I=0$,
$\frac13\sum_j(\fd_{I_j}+1+c_j)\log p$; that note reads the possible images maximally and takes a larger container, so
its bound exceeds the label average of the exact value of Theorem~\ref{t:exact} by
$\frac13\sum_j(1+c_j+\ceil{\lambda_j}-\lambda_j)\log p$. Its argument uses only the upper bound, so this does not affect
its conclusion.
\end{itemize}

\section{Relation to earlier work}\label{s:prior}
\begin{itemize}
\item \tagS\ The use of Proposition~1.1. \cite[Prop.~1.1]{IUT4} is stated with a free choice of $*$, while the proof of
Proposition~1.2(ii) uses $p^{\fd_I+a_I}$ (Remark~\ref{r:source}). Yamashita restates the $*$-form in his Lemma~1.4 and
uses $\ceil{\fd_I}$ in Lemma~1.6, obtaining $(-\lambda+\fd_I+4)\log p$ \cite[pp.~13--19]{Yam}. \tagX\ Dupuy and Hilado
write ``We can chose some $K_i$ which makes the inclusion tightest'' \cite[Rem.~2.8.2, p.~10]{DHPS}, and in step (6.4)
exempt the factor with the largest different \cite[p.~20]{DHPS}. Mochizuki--Fesenko--Hoshi--Minamide--Porowski change the
numerical constants of Theorem~1.10 but not Propositions~1.1--1.4 \cite[proof of Thm.~5.1]{MFHMP}.
\item \tagX\ The case $\lambda=0$. lana-agents/iut states (file \lean{Iut/WallEquiv/VcI.md}, first committed
13 October 2026): ``There the hull is IUT IV's container $p^{-d_I}\cdot(R_I)^{\sim}$ exactly (IUT IV, Prop. 1.4(iii)
with $\lambda = 0$, $I^* = \emptyset$; IUT's bound $(d_I + 1)\cdot\log p$ adds 1 for the rounding
$\lfloor\lambda - d_I - a_I\rfloor$).'' \cite{lana}. The author's \cite{Gap} computes the span of the orbit for the
log-shell-based region (Remark~\ref{r:frames}), using the same lattice lemma as the lower bound of
Theorem~\ref{t:exact}; that lemma is elementary and not new.
\item \tagX\ Bounded room. Mochizuki writes that ``automorphisms of the $\Zp$-module $I_k$ only give rise to bounded
discrepancies in the orders of nonzero elements of $\cO_k$'' \cite[Ex.~3.5.1]{EssLgc}, and that the confusion between
$q$-parameters and their large powers ``can never occur as a consequence of (Ind1), (Ind2), i.e., both of which amount
to automorphisms of the [underlying topological module of the] log-shells involved'' \cite[p.~101]{EssLgc}. In the proof
of Corollary~3.12, Step~(x), the resulting log-volumes ``are invariant with respect to the indeterminacies (Ind1),
(Ind2)'' \cite[p.~181]{IUT3}; this concerns each image, not the hull of their union. Takkun-kohinata/IUT\_LEAN
(15 July 2026) states, in a comment of \lean{Thm311.lean}, that (Ind1), (Ind2) preserve volume and so do not contribute
to the volume estimate of the hull (our paraphrase) \cite{Takkun}; again this concerns each image. Zeng computes the
index by which the hull enlarges the integral structure of a tensor product of rings of integers \cite{Zeng}.
\item \tagX\ Non-geometric automorphisms. Mochizuki writes that the valuation data ``is precisely the data that is
`deformed' by automorphisms of $k$ that do not arise from field automorphisms'' \cite[\S I3]{AbsTopIII}, describes (Ind1)
as arising from automorphisms of $G_k$ that do not in general preserve the ring structure \cite[\S3.7(i)]{Alien}, and in
an e-mail quoted by Joshi writes that such automorphisms ``can act in a fashion that fails to preserve differents,
discriminants'' and other quantities \cite[Appendix 16]{JoshiAnab}. Scholze and Stix note: ``Such isomorphisms exist for
Galois groups of $p$-adic fields, but this did not seem to enter the discussion in a critical way'' \cite[fn.~4]{SS}.
Mochizuki's formalization report mentions ``a small discrepancy (arising from (Ind1,2,3) + hull)'' \cite[p.~11]{Fm26}.
\item \tagX\ The classes. Dupuy and Hilado treat (Ind1), (Ind2) as acting through the automorphisms of the
$\Zp$-lattices of the log-shells \cite[p.~17, \S4.8]{DH1}. lana-agents/iut uses \texttt{tensorAut}, which contains the
image of $\Ism$ and is described as a deliberate enlargement \cite{lana}. Minamide and Tsujimura write ``Theorem F
implies that (Ind2) is relatively small'' \cite{MT}. Joshi reads (Ind1) as anabelomorphy and (Ind2) as a choice of
arithmetic holomorphic structure \cite[\S8.11.2]{JoshiIII}, a reading different from the one used here.
\item \tagP\ What is added here: the exact value with the ceiling for the region $p^{\lambda}\cdot\tI$ of Step~(v); the
difference of exactly $\log p$ from the container when $\lambda\notin\Z$, traced to the use of Proposition~1.1; the
arrangement of the four classes, with the non-geometric automorphisms of $G_v$ in the class $\cA$; and the reduction of
condition (U) to one condition on the transvections.
\end{itemize}

\section{Limitations}
\begin{itemize}
\item Only small places. $p=2$ and primes with $e_v>p-2$ (for instance the bad place above $59$ of the datum D211 of
\cite{Cat}, with $e=105$) are not treated.
\item The family of possible images omits (Ind3), takes the Kummer isomorphisms to be the identity, and treats label
permutations as Step~(v) does.
\item Condition (U) is not proved for $e_v\ge2$, and nothing is said for $\lambda\notin\Z$ about the actual class.
\item The effect on the final constants of \cite[Thm.~1.10]{IUT4} is not computed.
\end{itemize}

\section{Reproducibility and formalization}
The folder \texttt{three-tiers/} of the repository \url{https://github.com/turahigashi/iut-explicit-initial-theta}
(tag \texttt{v0.25.0}), also deposited with this note, contains the scripts and their outputs:
\texttt{check\_example.py} and \texttt{check\_nfactor.py} (the largest $k$ and the random factorwise images of
Remark~\ref{r:checks}), \texttt{check\_perfactor.py} (the same for two factors with 400 samples), and
\texttt{example\_d109.py} (the table of Section~\ref{s:example}, reading the place data of D109 from the program of
\cite{Walk}). They use exact rational arithmetic in Python.

Versions of Theorems~\ref{t:N}, \ref{t:Aupper} and \ref{t:Areach} and of Proposition~\ref{p:transvection} (with a
common ramification index above $p$ where the statements here allow different ones), and of the necessary condition in
the last remark of Section~\ref{s:actual}, were also checked in Lean~4 with Mathlib, for a model of the possible images
without (Ind3), in a development of the author that is not part of this deposit. The proofs given here do not rely on
it. Theorem~\ref{t:exact}, Corollary~\ref{c:compare} and Proposition~\ref{p:inside} are proved on paper only.

\section*{Disclosure}
The comparison studied here arose in the author's discussions with generative AI tools, and the author decided to make
it the subject of this note. The proofs were found with generative AI tools (large language models) under the author's
direction, and these tools were also used for the computations, the formalization, the literature search and the
writing. The author's own reading of the proofs is still in progress. The author takes full responsibility for the
note.

\begin{thebibliography}{AbsTopIII}\raggedright
\bibitem[AbsTopIII]{AbsTopIII}
Shinichi Mochizuki, \emph{Topics in absolute anabelian geometry III: global reconstruction algorithms}, version of
November 2015.

\bibitem[Alien]{Alien}
Shinichi Mochizuki, \emph{The mathematics of mutually alien copies: from Gaussian integrals to inter-universal
Teichm\"uller theory}, RIMS K\^oky\^uroku Bessatsu, preprint version.

\bibitem[Cat]{Cat}
Toshihisa Urahigashi, \emph{Explicit initial $\Theta$-data and the two readings of Step (v) of IUT IV, Theorem 1.10:
a catalogue}, version 1.2.0, Zenodo, 2026, \href{https://doi.org/10.5281/zenodo.24039374}{doi:10.5281/zenodo.24039374}.

\bibitem[DH1]{DH1}
Taylor Dupuy and Anton Hilado, \emph{The statement of Mochizuki's Corollary 3.12, initial theta data, and the first two
indeterminacies}, arXiv:2004.13228v1, 28 April 2020.

\bibitem[DHPS]{DHPS}
Taylor Dupuy and Anton Hilado, \emph{Probabilistic Szpiro, baby Szpiro, and explicit Szpiro from Mochizuki's
Corollary 3.12}, arXiv:2004.13108v2, 29 April 2020.

\bibitem[EssLgc]{EssLgc}
Shinichi Mochizuki, \emph{On the essential logical structure of inter-universal Teichm\"uller theory in terms of
logical AND ``$\wedge$''/logical OR ``$\vee$'' relations: report on the occasion of the publication of the four main
papers on inter-universal Teichm\"uller theory}, March 2024.

\bibitem[Fm26]{Fm26}
Shinichi Mochizuki, \emph{On the formalization of IUT: a preliminary progress report} [joint work in progress with
Y.~Hoshi, G.~Yamashita, Y.~Yang, \dots], April 2026.

\bibitem[Fourth]{Fourth}
Toshihisa Urahigashi, \emph{Split bad primes and the tuple structure of inter-universal Teichm\"uller theory: what the
source states, and what follows (a fourth report)}, version 1.0.5, Zenodo, 2026,
\href{https://doi.org/10.5281/zenodo.24039383}{doi:10.5281/zenodo.24039383}.

\bibitem[Gap]{Gap}
Toshihisa Urahigashi, \emph{The gap between the two walls is positive: a closed form for the (Ind2)-terms of
lana-agents/iut}, version 1.1.0, Zenodo, 2026,
\href{https://doi.org/10.5281/zenodo.24039377}{doi:10.5281/zenodo.24039377}.

\bibitem[HN]{HN}
Yuichiro Hoshi and Yu Nishio, \emph{On the outer automorphism groups of the absolute Galois groups of
mixed-characteristic local fields}, RIMS Preprint 1931, November 2020.

\bibitem[Hoshi]{Hoshi}
Yuichiro Hoshi, \emph{Introduction to mono-anabelian geometry}, RIMS Preprint 1868, January 2017.

\bibitem[HoshiTop]{HoshiTop}
Yuichiro Hoshi, \emph{Topics in the anabelian geometry of mixed-characteristic local fields}, Hiroshima Math. J.
\textbf{49} (2019), 323--398.

\bibitem[IUT3]{IUT3}
Shinichi Mochizuki, \emph{Inter-universal Teichm\"uller theory III: canonical splittings of the log-theta-lattice},
Publ. Res. Inst. Math. Sci. \textbf{57} (2021), 403--626.

\bibitem[IUT4]{IUT4}
Shinichi Mochizuki, \emph{Inter-universal Teichm\"uller theory IV: log-volume computations and set-theoretic
foundations}, Publ. Res. Inst. Math. Sci. \textbf{57} (2021), 627--723.

\bibitem[JoshiAnab]{JoshiAnab}
Kirti Joshi, \emph{On Mochizuki's idea of anabelomorphy and its applications}, arXiv:2003.01890v7, 21 June 2026.

\bibitem[JoshiIII]{JoshiIII}
Kirti Joshi, \emph{Construction of arithmetic Teichm\"uller spaces III: a `Rosetta Stone' and a proof of Mochizuki's
Corollary 3.12}, arXiv:2401.13508v4, February 2025.

\bibitem[JW]{JW}
Uwe Jannsen and Kay Wingberg, \emph{Die Struktur der absoluten Galoisgruppe $\mathfrak p$-adischer Zahlk\"orper},
Invent. Math. \textbf{70} (1982), 71--98.

\bibitem[Kondo]{Kondo}
Kaiji Kondo, \emph{Anabelian aspects of the outer automorphism groups of the absolute Galois groups of
mixed-characteristic local fields}, arXiv:2512.09231v2, 12 December 2025.

\bibitem[lana]{lana}
Project LANA, \emph{lana-agents/iut}, Lean repository, \url{https://github.com/lana-agents/iut}, commit
\texttt{553076a} (3 November 2026). Files cited: \lean{Iut/WallEquiv/VcI.md} (first committed 13 October 2026),
\lean{Iut/Concrete/LocalConstruct/IndAut.lean} (11 October 2026), \lean{Iut/Concrete/Ind1Invariance.lean}.

\bibitem[MFHMP]{MFHMP}
Shinichi Mochizuki, Ivan Fesenko, Yuichiro Hoshi, Arata Minamide and Wojciech Porowski, \emph{Explicit estimates in
inter-universal Teichm\"uller theory}, Kodai Math. J. \textbf{45} (2022).

\bibitem[MT]{MT}
Arata Minamide and Shota Tsujimura, \emph{Anabelian geometry for Henselian discrete valuation fields with quasi-finite
residues}, RIMS Preprint 1973, June 2023.

\bibitem[SS]{SS}
Peter Scholze and Jakob Stix, \emph{Why abc is still a conjecture}, 2018.

\bibitem[Takkun]{Takkun}
Takkun-kohinata, \emph{IUT\_LEAN}, Lean repository, \url{https://github.com/Takkun-kohinata/IUT_LEAN}, 15 July 2026.

\bibitem[Third]{Third}
Toshihisa Urahigashi, \emph{The cost of each distinction in the local computations of inter-universal Teichm\"uller
theory IV (a third report)}, version 1.0.3, Zenodo, 2026,
\href{https://doi.org/10.5281/zenodo.24039382}{doi:10.5281/zenodo.24039382}.

\bibitem[Walk]{Walk}
Toshihisa Urahigashi, \emph{A programmer's walkthrough of inter-universal Teichm\"uller theory}, version 1.2.0,
Zenodo, 2026, \href{https://doi.org/10.5281/zenodo.24039375}{doi:10.5281/zenodo.24039375}.

\bibitem[W311]{W311}
Toshihisa Urahigashi, \emph{Where Theorem 3.11 sits: the log-volume content of the results cited in the proof of
IUT III, Corollary 3.12, does not imply its conclusion}, version 1.1.0, Zenodo, 2026,
\href{https://doi.org/10.5281/zenodo.24039357}{doi:10.5281/zenodo.24039357}.

\bibitem[Yam]{Yam}
Go Yamashita, \emph{A proof of abc conjecture after Mochizuki}, survey, 31 August 2017.

\bibitem[Zeng]{Zeng}
Chunping Zeng, \emph{An integral-structure inflation in the hull operation of IUT Corollary 3.12}, Zenodo, 2026,
first version 29 October 2026, \href{https://doi.org/10.5281/zenodo.23754276}{doi:10.5281/zenodo.23754276}.
\end{thebibliography}

\end{document}
